\begin{CJK}{UTF8}{gbsn}
\chapter{摘要}
	\section{引言}\label{intro_chinese}
	论文基于CMS探测器收集的 质心系能量为 =13 TeV 的质子-质子对撞数据，在 $ H \to 4\ell$ ($\ell = \Pe$, $\mu$) 衰变道测量Higgs粒子的基准截面 (积分和微分)。测量中使用Higgs质量 以及大型强子对撞机（LHC）产生的积分亮度为 （Full Run II）的实验数据。其中微分截面被表示为几个动力学变量的函数，这些动力学变量依赖于Higgs粒子的产生机制，该测量同时也在TeV能标上直接测试了微扰QCD计算。
	\subsection{意义}\label{motivation_chinese}
不依赖于理论模型，给定过程的产生截面是Higgs玻色子的重要性质，此处产生截面给出该给定过程发生的可能性。这篇论文，通过CMS收集的　质心系能量 =13TeV 的质子-质子对撞数据，在　$ H \to 4\ell$ decays ($\ell = \Pe$, $\mu$) 衰变道测量Higgs粒子的积分，微分基准产生截面．考虑到以下因素该论文基于Higgs动力学变量测量微分截面。Higgs横向动量可以用于检测 pQCD的计算（特别是当实验中使用全举多jet 末态），另外这个变量对Higgs的拉适量的结构非常敏感，Higgs的快度对该实验中的部分子分布函数（PDFs）非常敏感。jet相关的观测量对理论模型非常敏感，boosted夸克，胶子辐射也在这篇论文中讨论。目前的测量结果与标准模型预言相一致，但是，在此类测量结果中，任何与标准模型的偏差，都直接指向超越标准模型的贡献。Ref. ~\cite{Khachatryan:2015yvw, CMS:2015hja, CMS:2016rqf, Sirunyan:2017exp, CMS:2018hhg, ATLAS-CONF-2019-029, CMSHggFiducial8TeV, atlas-conf-2019-025, ATLAS:2018uso, ATLAS-CONF-2019-032, ATLAS:2017myr,  CMS:Hgg_8TeV, CMS:2019chr, Sirunyan:2020xwk} 是CMS，ATLAS合作组在不同Higgs衰变道中对这些问题的研究。
	\section{理论背景}\label{theory_motivation_chinese}
此次测量是为精确检验粒子物理标准模型（SM），经过长时间的研究，物质及物质间的相互作用被标准模型精确描述（电磁力，弱力，强力不包括引力）。这些基本作用力的细节描述在\ref{fundamentalforces}。Higgs玻色子时粒子物理标准模型的重要组成部分，它被在LHC上运行的质子-质子实验所首先发现。标准模型对基本相互作用力的描述由数学的对称法则组成。它提供的数学描述只能适用于质量为0 的粒子。Higgs机制的引入，使得标准模型更加普世，解决了标准模型只能描述0质量粒子的问题。在Higgs机制中，粒子通过电弱对称性破缺获得质量。Higgs粒子的存在就是Higgs机制的直接证明 ~\cite{Reference7,Reference8,Reference9,Reference10,Reference11,Reference12}。
标准模型是用 \suthreec$\otimes$\sutwol$\otimes$\uoneg~ 规范对称群描述的量子场论。它的公式化源自不同时期的不同科学家跨越第20世纪的相互协作。1961年，Glashow提出电磁力和弱力有相似的结构，可以统一用 \sutwol$\otimes$\uoneg ~对称群解释。Brout, Englert和 Higgs 提出无质量的基本粒子可以通过自发对称性破缺以及相关的另一个基本粒子Higgs 玻色子获得质量。之后，Salam和Weinberg提出Glashow的 \sutwol$\otimes$\uoneg 模型中L表示左手征态费米子场，Y是弱超荷与弱同位旋和电荷相关。到这里完成了标准模型中的电弱统一。如 \ref{ewk_glashow} 中的描述，Glashow对电磁力和弱力的统一描述只适用于无质量粒子。然而规范场必须要求粒子拥有质量，以描述若相互作用以及和实验结果相一致。Higgs机制被提出作为标准模型的一部分，同时引入一个新的标量场，同时这个标量场的变换满足规范协变性和重整化。自发对称破缺的公式化概念以及费米子和玻色子（W，Z）的质量起源详见 ~\cite{Reference7,Reference8,Reference9,Reference10,Reference11,Reference12}关于Higgs粒子的细节描述在\ref{chapter1_sec2}。
	\subsection{LHC中质子-质子对撞现象学} \label{lhc_pheno_chinese}
	LHC是同类加速器物理实验中最大的实验装置，建造它的目的是为了精确检验粒子物理标准模型。质子由三个价夸克（uud）通过交换胶子的强相互作用束缚在一起。夸克通过一些方法在质子内部相互作用。一个夸克与另外两个夸克彼此相互作用，或者通过湮灭过程产生胶子，胶子进一步产生夸克反夸克对。结果导致除了价夸克外，质子内部还有海夸克和胶子（在一起被称为部分子）。在对撞实验中，大的动量转自可能发生在对撞的部分子之间。这被称为硬散射，一些共振态可能在这个过程中产生比如Higgs。一个过程发生的可能性被定义为该过程的截面。LHC中质子-质子对撞的不同过程的产生截面见图 \ref{Fig:lhc_all}。没有参与硬散射的部分子通常会直接向前飞走后这参与低动量过程。低动量过程事例被称为underlying 事例。如果在个质子-质子对撞事例中有超过一个相互作用发生，成为多部分子相互作用（MPI）。
被加速的带电部分子可能会向外辐射光子和胶子，被称为轫致辐射。如果这些辐射可以关联到一个入射部分子，这辐射成为初态辐射。如果关联到一个出射部分子，成为末态辐射。之后夸克和胶子通过强子化过程（相互束缚形成色单态）组成一些强子。这些强子经产表现为喷注，称为jet，如\ref{sec:jets}中的描述。
在硬散射中，表示一个过程发生可能性的截面可以通过Factorization ($\mu_{F}$) 定理计算。这个定理使用Factorization因子帮助我们将截面分解为硬，软过程。过于截面的数学描述相见 \ref{phenomenology}。

	\subsection{}\label{higgs_prod_dec_chinese}
该论文中主要设计5个不同的Higgs产生模式。在LHC中最主要的Higgs产生模式是胶子-胶子融合。胶子通过虚夸克圈图与Higgs耦合。因为耦合直接与费米子质量成正比，所i有top夸克圈图在这里占主导。在LHC和LEP实验中VBF是第二重要的Higgs产生模式。在VBF产生模式中，两个W或者Z 玻色子融合产生Higgs粒子。VBF过程的实验信号因为带有两个向前的高能jet的特征，所以可以与胶子融合信号有很好的区分。同时这个特征也用来压低背景。在W/Z联合产生模式中，Higgs粒子由高能W或Z玻色子产生，这个W或Z玻色子由两个费米子融合产生。该过程在LEP实验中是最主要的Higgs产生过程，Tevatron实验中是次级产生过程。在LEP电子正电子实验中，该过程产生了很多W和Z玻色子。在强子对撞机中，Tevatron使用质子反质子对撞，这与LHC更相似。最小的Higgs产生过程如图~\ref{Fig:higgs_production} 所示。该过程由两个胶子碰撞产生top和反-top夸克对，之后$t\bar{t}$融合产生Higgs。
标准模型126 GeV Higgs粒子的寿命在 $1.6\times10^{-22}$~s  它可以衰变向所有标准模型有质量粒子（在壳重粒子的可能性更大）。Higgs玻色子向不同粒子衰变的分支比如图~\ref{Fig:Higgs_BR}所示 ~\cite{Referenceb}。$H\rightarrow b\bar{b}$（$m_H \approx 125$ ~GeV）是最重要的衰变模式。但是该衰变道的QCD背景比信号的高几个数量级，所以该衰变道的零密度最低。然而，CMS和ATLAS合作组以及一些其他实验比如 ~\cite{Reference13}, ~\cite{Reference14} 和 ~\cite{Reference15}（CDF和D0）已经利用了这个衰变道。B夸克在探测器中也很难重建。当它们的在探测器中衰变时，它们的碎片产生很多jet和强子。我们很难完美的重建这些jet。正是因为如此，我们必须开发其他衰变道。Higgs衰变到双光子是非常重要的衰变道，因为它有两个标志性的高能光子和好的质量分辨率。
因为Higgs玻色子的质量在125GeV，Higgs玻色子也可以衰变到两个矢量玻色子。这里摸一个或者两个玻色子可能是不在壳的。WW衰变可以进一步衰变到中微子，但是中微子在探测器中不可能完全重建。$H \rightarrow ZZ \rightarrow \ell\ell\ell\ell$ 衰变道被视为黄金衰变道，尽管衰变分支比很低。这些轻子是高能轻子，可以在探测器对它们的重建质量很高~\cite{Referenceb}。改论文中对Higgs粒子性质的测量正是基于$H \rightarrow ZZ \rightarrow \ell\ell\ell\ell$衰变道。

	\section{实验背景} \label{exp_bkg_chinese}
	在Tevatron，LEP高能物理实验的精确测量中，已经证明标准模型的正确性，但是仍然没有观测到Higgs粒子。在2010年，大型强子对撞机被建造，用来对标准模型做更多的检验，以及验证之前的实验结果。LHC安装在CERN（欧洲核子研究中心）坐落在瑞士-法国边境，日内瓦近郊。它是在粒子物理同类实验中，能量最高的实验装置。它有足够的能力将质子束流加速到14TeV 亮度峰值为$\mathscr{L}$ = 10$^{34}$cm$^{2}$s$^{-1}$。它是人类有史以来建造过的最先进，最复杂和最大的实验装置，由来自超过100个国家，10000名工程技术人员以及科学家共同设计建造。从1998年到2008年，LHC被建造在地下75-150米，3.8米宽，27千米长的隧道里，该隧道在1983-1988年曾经用来安装LEP对撞机。在LHC中两个反向旋转的质子束流运行在高真空的管道中以避免束流和空气分子的额外碰撞。2012年，在LHC上运行的两个实验CMS和ATLAS联合发布发现了一个新的重标量玻色子，Higgs玻色子 ~\cite{Reference20,Reference21}。基于这个发现，物理学家希望通过精确测量Higgs粒子的性质，以进一步检验标准模型。Higgs粒子的质量是Higgs机制中的唯一自由参数，它的测量结果是 $m_H = 125.09\pm 0.21(stat.)\pm0.11(syst.)$~GeV~\cite{Reference22}。 
在运行第一阶段成功取数后，大型强子对撞机LHC关闭了两年。在2013年至2014年间LHC设备和探测器进行了重大升级。在2015年重启了LHC，质心系能量达到13TeV（接近其设计值14TeV）。如图所示，在2015年LHC累积了质心系能量13TeV积分亮度为4.2fb-1的数据。随后在2016，2017，2018年间，分别累计了41.0fb-1，49.8fb-1和67.9fb-1的数据，在第二阶段运行后，LHC目前正在经历第二阶段的长期关闭。在此期间，几处升级改造正在进行，为后期LHC高亮度运行（HL-LHC）做准备。但是，这方面的主要升级将在第三次关闭期间进行（即大型强子对撞机运行第3阶段之后）。其一是用Linac4替换Linac2。这个长度约为90米，耗时约10年建成，用于加速负氢离子（H-），使粒子束能量达到160 MeV，并使其强度翻倍。此外，在第二阶段下线期间，LHC上各个实验组，如CMS，ATLAS，ALICE和LHCb，也计划进行相应的升级改造。
ATLAS正在构建与触发系统集成的缪子小轮。为了应对高堆积（pileup），LAr电子设备也正在升级。
CMS将强子量能器桶部和端盖的光电探测器（HPD）更改为硅光电倍增管（SiPM）。此外，用铝代替不锈钢制作束管。由于辐射造成了强子量能器端盖的损伤， 
闪烁条也将更换。为了准备高积分亮度运行（HL-LHC），CSC缪子室正在升级改造，装备新的前端末端（FE）电子学器件。
ALICE正在升级内部径迹系统（ITS），时间投影室（TPS），中央触发处理器，数据采集（DAQ）高级别触发（HLT），缪子前向径迹探测器（MFT）等。
LHCb将其读数从1 MHz提高到40 MHz，将用完全基于软件的触发方式替代前端末端（FE）电子学器件以及使用全新的计算基础架构。
	\subsection{紧凑型缪子硅探测器}\label{cms_chinese}
	本文使用的数据来于CMS探测器。有关CMS探测器的构造、内部结构以及CMS检测器如何重建事件的说明，请参见章节 \ref{Chapter2_CMS}。大型强子对撞机为粒子物理实验研究开辟了新领域，基于最新技术，配备最先进的实验设施。为了确保实验结果的有效性，LHC 有两个主要的目标基本一致的实验进行交叉验证，cms是其中之一。它们具有相同的物理目标，但探测器的设计和技术不同。CMS能够在广范围内为μ识别和动量提供良好的分辨率，其动量在1 TeV上仍有极致的精确度。CMS探测器具有良好的带电粒子内部径迹重建效率和分辨率。它具有高效的触发系统，并能更好地识别τ和b喷注。此外，双电子、双光子、双喷注以及缺失横向能量（MET）具有优异的质量和能量分辨率。这使其能够研究标准物理（包括希格斯物理），额外维度和暗物质，覆盖100GeV到TeV能量尺度的物理~\cite{Reference34}。作为大型强子对撞机上两个广物理目标探测器之一，CMS探测器是圆柱形，长度约为28 m，外半径为大约7 m，总高度约为15 m，总重量约为14000吨。CMS检测器的整体设计如图??所示。依据CERN早期的实验经验，cms建造不是全在地下完成。cms探测器被分为十五份，然后放置地表下在LHC隧道中组建。
\par 从对撞点径向向外看分别是CMS的内部径迹探测器（ITS），电磁量能器（ECAL）和强子量能器（HCAL）和缪子室，这些子探测器处在大型磁场内。CMS的内部径迹探测器具有较大的径迹多样性，因为它最接近发生对撞的束管。电磁量能器由钨酸铅晶体构成，涵盖了赝快度小于3.0的范围。在电磁量能器后是由黄铜制成的强子量能器。CMS探测器的最外层是缪子系统。
\par LHC在每秒发生4000万次碰撞（40 MHz）。但是并不是所有的碰撞都是感兴趣的研究物理目标。此外，数据采集系统（DAQ）不能有效地高速率读取每个事件。存储碰撞中产生的无趣事件也将消耗大量的存储。因此，需要触发系统，用以减少在碰撞期间记录事件的速率。触发系统将选择出仅对物理目标至关重要的事件。CMS的触发系统是一个复杂的两级系统，包括第一层触发（L1）和高级别触发（HLT）。
	\section{希格斯基准截面的测量策略}\label{strategy_chinese}
	标准模型的希格斯玻色子有五个主要的生产模式（ggH，VBF，WH，ZH和ttH），由{\sc powheg V2}产生子产生这些信号样本。WH和ZH的生产通过{\sc powheg}的{\sc MiNLO HVJ}扩展进行模拟~\cite{Luisoni:2013kna}。希格斯到四轻子衰变过程由{\sc JHUGEN}产生子模拟\cite{Gao:2010qx}。产生WH，ZH和ttH样本时，允许希格斯衰变到H$\to ZZ \to 2\ell2$X，相应地，四轻子事件的两个轻子有可能分别从关联的W，Z玻色子或顶夸克的衰变中产生。使用{\sc pythia8.209}进行部分子的簇射，并且在所有情况下，簇射的匹配标准是先允许所有能量范围的量子色动力学QCD发散，然后在POWHEG等级下进行判选。所有样本均使用默认的NNPDF 3.1 NLO 部分子分布函数（PDF）~\cite{Ball:2014uwa}。这些信号样本截面信息包含在 ~\cite{CMS:2019chr}。这些不同生产模式的样本还用于评估测量结果对模型的依赖性，如表~\ref{tab:efficiencySM}所示。关于评估模型依赖性的方法的详细信息将在\ref{sec:Systematics}节中进行描述。ggH生产模型的{\sc powheg}和{\sc NNLOPS}样本被用以标准模型计算的理论预测，以及与测量结果进行比较。本底估计中使用{\sc powheg V2}和{\sc pythia8}交互产生次领头阶（NLO）精度的主导过程$\Pq\Paq \to \Z \Z \to 4\ell$。$\Pg\Pg \! \to \! ZZ$过程使用{\sc MCFM}进行领头阶精度（LO）的模拟。有关模拟和数据样本的更多详细信息，请参见\ref{sec:Datasets} and ~\ref{sec:Simulations}部分。
希格斯截面的测量，包括积分截面和微分截面，在基准相空间中进行，重建选择与基准空间进行匹配。对于希格斯的寻找，首先从选定的轻子中构建四轻子候选者，轻子的选择条件包括顶点约束${\rm SIP_{3D}} < 4$~ (${\rm SIP_{3D}}$是个重要的影响参数，其值是将影响参数除以相关的误差），以及减去其中末态辐射光子贡献的孤立性选择（isolation cut）。Z玻色子候选者由相反电荷的相同风味的轻子($e^+ e^-$, \, $\mu^+\mu^-$)构成，并且要求其不变质量在12GeV到120GeV范围内。然后从不重叠的Z玻色子候选对象构建ZZ候选对象。其中不变质量更接近于标准模型Z玻色子的被标记为${\rm Z_1}$，另一个为${\rm Z_2}$。信号区间由至少要包含一个ZZ候选者的事例构成，更多细节参见章节\ref{sec:events_selection}。
\par $pp\to$H$\to4\ell$的积分和微分基准截面的测量由信号和本底参数化的最大似然拟合得到，包含四轻子质量分布，观测量的个数，并直接从中提取基准截面。系统误差以干扰参数的形式包括在内，并在拟合过程中得到有效整合。使用渐近方法~\cite{LHC-HCG}和基于轮廓似然比~\cite{Cowan:2010js}的检验统计量获得结果。从观察到的分布中展开探测器效应的过程与参考文献~\cite{CMSH4lFiducial8TeV}和~\cite{CMSHggFiducial8TeV}中的相同。在微分截面测量中，将有限的效率和分辨率效果编码到探测器响应矩阵中，该矩阵描述事件如何从基准级别的既定观察量单元映射到重建级别的既定单元。该矩阵是对角线主导的，较多的非对角线元素用于涉及喷注的观测量。
\par 系统误差的主要来源是实验误差，包含轻子鉴别效率和积分亮度测量，理论误差来源被认为是次重要的。为了评估测量的模型依赖性，使用不同的响应矩阵重复展开过程，这些响应矩阵是通过在其实验约束范围内更改每个标准模型产生模式的相对部分而创建的。相对于实验系统误差，这个不确定性认为可以忽略。

	\section{结果和展望} \label{results_chinese}
	为了使所有测量结果独立于希格斯的产生方式，基于轻子运动学和事件拓扑结构定义基准相空间，所有的截面测量都取自基准相空间。在质心系能量13TeV下，测量的积分基准截面为$2.73^{+0.30}_{-0.29}=2.73^{+0.23}_{-0.22}({\rm stat.})^{+0.24}_{-0.19}({\rm syst.})~{\rm fb}$。基准微分截面表述为快度、四轻子体系的横动量、相关的多重喷注以及领头喷注的横动量的函数。微分截面在这些观测量的各个单元内进行测量。对于每个运动学或微分观测量，优化每个单元边界的选择过程，以使观测量的每个单元具有相似的预期测量横截面。微分横截面测量应在观测量的每个单元上的重建层次上减去背景。对于所有微分单元而言，本底的质量分布是不相同，因为它没有共振结构。对于不可区分和可区分的背景，分别通过模拟事例和数据中的控制区域，通过四轻子质量模板形状获得了四轻子质量谱，分别用于观测量的微分测量的每个单元，如参考文献中所述~\cite{Chatrchyan:2013mxa}。
	\par 这些测量结果对于检验标准模型希格斯部分的预测至关重要，并且可以用于许多超标准模型理论的相空间约束。将所有积分和微分测量值与基于标准模型的理论预测值进行比较，发现在误差范围内是一致的。需要更精确的希格斯玻色子性质测量来发现与标准模型预测的任何可能偏差，这可能是通往新物理学的大门。
	\par 这些事情是可能的，但是超出了本文的范围。当累积了更多的数据，我们可以获得更精确的结果。使用全部第二运行阶段的数据，研究了更为细致的微分截面事，涉及了更多的观测量。当前测量是基于{\sc powheg}和{\sc NNLOPS}的预测。但是，可以使用更多的理论预测，用于与微分单元中观察到的横截面进行比较。此外，可以研究微分观测量的潜在理论解释。

\end{CJK}
